\section{Research Plan and Future Timeline}
\label{sec:research-plan}

% Question: What remains to be done, and in what order?
% Answer: Five work packages, from closing the WASP-17 b benchmark to a
% cross-instrument synthesis and the thesis; Figure 5 gives the timeline.
% Evidence: The previous, longer plan (git history of this file); the Methods.
The remainder of the PhD will turn the current development system into a
tested method, determine where reduction choices change atmospheric
interpretation, and then combine validated observations across visits,
instruments and observing geometries. The programme is organised into five
work packages, shown in Figure~\ref{fig:gantt}; the timeline begins with the
remaining part of 2026~Q3 and contains future work only.

\paragraph{WP1: NOVA and injector closure.}
The first priority, to be completed by 30 September 2026, is to close the
WASP-17~b benchmark: the pilot recoveries and the envelope of
Section~\ref{sec:envelope}, with strictly converged NOVA fits. The benchmark will then be expanded to a
preregistered matrix which includes a featureless spectrum, broad molecular
structure, localised features, several signal amplitudes, several noise
realisations, different background amplitudes and morphologies, and a
controlled mismatch of the instrument response. The
deliverable is a versioned benchmark package, and its main scientific result
will be a map of recovery error as a function of the injected and detector
conditions, not a single favourable spectrum.

\paragraph{WP2: Robustness, pipeline benchmark and methods paper.}
NOVA will next be run on out-of-transit data from further visits, to test
whether its response and background model are robust to changes in pointing,
background and detector state. The comparison will be extended to Eureka! and
the official JWST pipeline with ATOCA extraction \citep{DarveauBernierEtAl2022},
and a separate diagnostic track will pass every pipeline's extracted products
to a common light-curve fitter, to separate extraction effects from
differences in transit modelling. Ablations on the same injected data will
fix the background instead of fitting it, simplify the response, remove the
robust weighting, change the spectral basis, fit the orders separately and
reproduce selected extraction-first compressions, so as to separate the
effect of joint detector modelling from that of a particular nuisance
model. The main deliverable of WP1 and WP2 is the NOVA and
benchmark methods paper, to be submitted in 2027~Q1 together with the
software needed to reproduce its figures.

\paragraph{WP3: Multi-target NIRISS/SOSS validation.}
Once it passes the benchmark, the method will be applied to further public
SOSS transmission observations spanning stellar brightness, background and signal amplitude,
with priority for those with at least two published reductions, and
validated on real data by detector-domain predictive checks, consistency
between the orders, independent reductions and injection into the observation
residuals. The question is whether
NOVA changes repeatable atmospheric features or mainly offsets and
uncertainties.
The SOSS analysis will be frozen by 30 June 2027. A target-specific paper will
only be written where the recovered spectrum changes repeatably and the change
alters an atmospheric inference, and such papers may continue into 2027~Q3.

\paragraph{WP4: Emission, phase curves and other instruments.}
The model will be extended to secondary eclipses and phase curves, first
within SOSS, so that a change of geometry is separated from a change of
instrument, and then to NIRSpec and MIRI, each with its own benchmark. Suitable same-planet data exist: WASP-17~b has SOSS
transmission and eclipse measurements and consecutive NIRSpec transit and
eclipse observations \citep{LouieEtAl2025,GressierEtAl2024,
LustigYaegerEtAl2025}, WASP-121~b has NIRSpec and NIRISS phase curves and a
NIRSpec transmission analysis \citep{MikalEvansEtAl2023,GappEtAl2025,
SplinterEtAl2025}, and GJ~1214~b has a MIRI phase curve and a panchromatic
transmission analysis \citep{KemptonEtAl2023,OhnoEtAl2025}. The final
objective is a joint retrieval of selected planets across instruments and
geometries, with reduction uncertainty, visit variability and shared
parameters represented explicitly, so that it can test whether the combined
data tighten the constraints without hiding inconsistencies.

\paragraph{WP5: Synthesis and thesis.}
Writing will proceed alongside the papers. The research phase ends on approximately 27 April 2029, after accounting for
a one-month interruption, and the formal writing-up period, which runs to
28 October 2029, is reserved for completing and submitting the thesis, the
corrections and viva preparation, not for new method development.

\begin{figure*}[t]
\centering
\resizebox{\textwidth}{!}{%
\begin{ganttchart}[
  hgrid,
  vgrid,
  x unit=0.95cm,
  y unit title=0.42cm,
  y unit chart=0.31cm,
  title/.style={fill=black!6,draw=black!25},
  title label font=\bfseries\scriptsize,
  bar label font=\scriptsize,
  group label font=\bfseries\scriptsize,
  milestone label font=\scriptsize,
  bar height=0.55,
  group height=0.34,
  group peaks height=0.16,
  milestone height=0.45
]{1}{14}
\gantttitle{2026}{2}\gantttitle{2027}{4}\gantttitle{2028}{4}\gantttitle{2029}{4}\\
\gantttitle{Q3}{1}\gantttitle{Q4}{1}
\gantttitle{Q1}{1}\gantttitle{Q2}{1}\gantttitle{Q3}{1}\gantttitle{Q4}{1}
\gantttitle{Q1}{1}\gantttitle{Q2}{1}\gantttitle{Q3}{1}\gantttitle{Q4}{1}
\gantttitle{Q1}{1}\gantttitle{Q2}{1}\gantttitle{Q3}{1}\gantttitle{Q4}{1}\\
\ganttgroup[group/.append style={fill=ganttsoss!35,draw=ganttsoss}]
  {SOSS method and benchmark}{1}{4}\\
\ganttbar[bar/.append style={fill=ganttsoss!70,draw=ganttsoss}]
  {Injector and strict NOVA}{1}{1}\\
\ganttbar[bar/.append style={fill=ganttsoss!70,draw=ganttsoss}]
  {Multi-visit OOT robustness}{1}{2}\\
\ganttbar[bar/.append style={fill=ganttsoss!70,draw=ganttsoss}]
  {Expanded pipeline benchmark}{1}{2}\\
\ganttbar[bar/.append style={fill=ganttsoss!70,draw=ganttsoss}]
  {NOVA and benchmark paper}{2}{3}\\
\ganttbar[bar/.append style={fill=ganttsoss!70,draw=ganttsoss}]
  {SOSS target survey}{3}{4}\\
\ganttbar[bar/.append style={fill=ganttsoss!70,draw=ganttsoss}]
  {Conditional target papers}{4}{5}\\
\ganttmilestone[milestone/.append style={fill=ganttsoss,draw=ganttsoss}]
  {SOSS analysis freeze}{4}\\
\ganttgroup[group/.append style={fill=ganttgeneral!35,draw=ganttgeneral}]
  {Geometry and instrument generalisation}{5}{10}\\
\ganttbar[bar/.append style={fill=ganttgeneral!70,draw=ganttgeneral}]
  {Emission and phase curves}{5}{7}\\
\ganttbar[bar/.append style={fill=ganttgeneral!70,draw=ganttgeneral}]
  {Same-planet data audit}{5}{6}\\
\ganttbar[bar/.append style={fill=ganttgeneral!70,draw=ganttgeneral}]
  {NIRSpec extension}{7}{9}\\
\ganttbar[bar/.append style={fill=ganttgeneral!70,draw=ganttgeneral}]
  {MIRI extension}{8}{10}\\
\ganttgroup[group/.append style={fill=ganttsynthesis!35,draw=ganttsynthesis}]
  {Joint atmospheric inference}{9}{12}\\
\ganttbar[bar/.append style={fill=ganttsynthesis!70,draw=ganttsynthesis}]
  {Cross-geometry joint retrieval}{9}{12}\\
\ganttbar[bar/.append style={fill=ganttsynthesis!70,draw=ganttsynthesis}]
  {Synthesis paper}{11}{12}\\
\ganttmilestone[milestone/.append style={fill=ganttsynthesis,draw=ganttsynthesis}]
  {Science freeze}{12}\\
\ganttgroup[group/.append style={fill=ganttwrite!35,draw=ganttwrite}]
  {Thesis}{5}{14}\\
\ganttbar[bar/.append style={fill=ganttwrite!70,draw=ganttwrite}]
  {Thesis integration}{5}{12}\\
\ganttbar[bar/.append style={fill=ganttwrite!70,draw=ganttwrite}]
  {Formal writing-up period}{12}{14}\\
\ganttmilestone[milestone/.append style={fill=ganttwrite,draw=ganttwrite}]
  {Thesis submission}{14}
\end{ganttchart}%
}
\caption{Future-only PhD plan in quarters. The first column covers the remaining
part of 2026 Q3. The SOSS analysis is frozen by 30 June 2027; conditional target
manuscripts may continue into Q3. The formal writing-up period runs from
approximately 27 April to 28 October 2029.}
\label{fig:gantt}
\end{figure*}

\paragraph{Dependencies and risks.}
The injector and strictly converged NOVA fits precede the headline
comparison, atmospheric claims are only made after the benchmark checks are
passed, and a milestone is complete only when its products, configuration,
diagnostics and interpretation are archived together. The main technical risk
is an instrument response too inaccurate for strict recovery, which I will
address with held-out validation, independent calibration routes and
simulations with a known response; if it persists, the thesis can still quantify
the accuracy required. If NOVA does not outperform the extraction-based analyses,
a fair test still establishes when extraction is sufficient, and if time is
short, the synthesis can focus on one planet and two instruments.
